\documentclass[11pt]{article}
\usepackage[a4paper,margin=1in]{geometry}
\usepackage[T1]{fontenc}
\usepackage[utf8]{inputenc}
\usepackage{lmodern,microtype}
\usepackage{amsmath,amssymb,amsthm,mathtools}
\usepackage{enumitem,xcolor,hyperref}
\usepackage[nameinlink,capitalise,noabbrev]{cleveref}
\hypersetup{colorlinks=true,linkcolor=blue!60!black,citecolor=blue!60!black,urlcolor=blue!60!black}

\newtheorem{definition}{Definition}[section]
\newtheorem{assumption}[definition]{Assumption}
\newtheorem{theorem}[definition]{Theorem}
\newtheorem{proposition}[definition]{Proposition}
\newtheorem{corollary}[definition]{Corollary}
\newtheorem{conjecture}[definition]{Conjecture}
\newtheorem{remark}[definition]{Remark}

\newcommand{\C}{\mathbb C}
\newcommand{\R}{\mathbb R}
\newcommand{\E}{\mathcal E}
\newcommand{\Q}{\mathcal Q}
\newcommand{\Oterm}{\mathcal O}
\newcommand{\norm}[1]{\left\lVert #1\right\rVert}
\DeclareMathOperator{\crit}{crit}
\DeclareMathOperator{\dist}{dist}

\title{Lojasiewicz Pandrosion Descent\\[3pt]
\large Analytic residual energies, stable finite-slope steps, and convergence under gated corrections}
\author{Ivan Besevic\\Independent Mathematical Research}
\date{May 2026}

\begin{document}
\maketitle

\begin{abstract}
Previous Pandrosion descent papers isolated conditional Lyapunov inequalities for stable finite-slope pairs and for gated higher-order corrections. This paper connects those inequalities to analytic convergence theory. Let \(G:U\subset\R^n\to\R^n\) or \(G:U\subset\C^n\to\C^n\) be analytic and set
\[
        \E(x)=\frac12\norm{G(x)}^2.
\]
A Pandrosion sequence is called Lojasiewicz-admissible if it is bounded, satisfies sufficient decrease, and satisfies a relative-error estimate linking \(\nabla \E\) to the accepted step. Under the classical Lojasiewicz gradient inequality near any cluster point, every admissible sequence has finite length and converges to a critical point of \(\E\). If the cluster point is a regular zero of \(G\), the sequence converges to that zero. We then show that stable finite-slope Pandrosion steps with line search satisfy the sufficient-decrease part, and that small gated finite-slope Halley corrections preserve admissibility. The resulting convergence theorem is intentionally conditional: it does not claim that every probe is good, but it gives a rigorous target for probe selection, atlas normalization, equal-value pole avoidance, and Halley gating.
\end{abstract}

\paragraph{Scope and status.}
This paper is a convergence bridge. Its theorems use standard analytic Lojasiewicz/Kurdyka-Lojasiewicz hypotheses stated explicitly, while the final global avoidance statement remains conjectural.

\section{Analytic energy and Pandrosion steps}

Let \(G:U\to\R^n\) be real analytic. Complex analytic systems may be read as real analytic systems on \(\R^{2n}\). Define the residual energy
\[
        \E(x)=\frac12\norm{G(x)}^2.
\]
A finite-slope Pandrosion step at \(x_k\) selects a probe \(b_k\), a finite-slope matrix \(\Q_k=\Q_G(x_k,b_k)\), and solves
\[
        \Q_k\delta_k=-G(x_k).
\]
The accepted iterate has the form
\[
        x_{k+1}=x_k+\lambda_k(\delta_k+e_k),
\]
where \(0<\lambda_k\le1\) is a line-search parameter and \(e_k\) is an optional gated correction, for example a finite-slope Halley correction.

\section{Admissible descent sequences}

\begin{definition}[Lojasiewicz-admissible sequence]
A bounded sequence \((x_k)\subset U\) is Lojasiewicz-admissible for \(\E\) if there exist constants \(\sigma>0\), \(L>0\), and \(k_0\) such that for all \(k\ge k_0\),
\begin{align}
        \E(x_{k+1})+\sigma\norm{x_{k+1}-x_k}^2 &\le \E(x_k), \label{eq:suff-decrease}\\
        \norm{\nabla \E(x_{k+1})} &\le L\norm{x_{k+1}-x_k}. \label{eq:relative-error}
\end{align}
\end{definition}

\begin{remark}
The sufficient-decrease inequality is the Lyapunov part. The relative-error inequality says that accepted steps are not arbitrary energy-lowering jumps; the step length controls stationarity of the next point. This is the standard structure behind many analytic descent convergence theorems.
\end{remark}

\begin{assumption}[Lojasiewicz gradient inequality]
At every cluster point \(x_\ast\) of the sequence, there are constants \(c>0\), \(\theta\in[0,1)\), and a neighborhood \(V\) of \(x_\ast\) such that
\[
        |\E(x)-\E(x_\ast)|^\theta
        \le
        c\norm{\nabla \E(x)}
        \qquad (x\in V).
\]
\end{assumption}

\section{Finite-length convergence}

\begin{theorem}[Lojasiewicz convergence of admissible Pandrosion sequences]\label{thm:kl}
Let \(G\) be analytic and let \((x_k)\) be a Lojasiewicz-admissible sequence for \(\E=\frac12\norm{G}^2\). If the sequence has a cluster point \(x_\ast\), then \((x_k)\) has finite length in a neighborhood of \(x_\ast\):
\[
        \sum_{k\ge k_1}\norm{x_{k+1}-x_k}<\infty
\]
for some \(k_1\). Consequently the whole tail converges to \(x_\ast\), and \(x_\ast\) is a critical point of \(\E\).
\end{theorem}

\begin{proof}
The decrease inequality implies that \(\E(x_k)\) is nonincreasing and converges to some value \(\E_\infty\). Since \(x_\ast\) is a cluster point and \(\E\) is continuous, \(\E_\infty=\E(x_\ast)\). For all large \(k\), the iterates lie in a Lojasiewicz neighborhood of \(x_\ast\).

Set \(\Delta_k=\norm{x_{k+1}-x_k}\). From \eqref{eq:suff-decrease},
\[
        \sigma \Delta_k^2\le \E(x_k)-\E(x_{k+1}).
\]
The Lojasiewicz inequality and the relative-error bound give
\[
        |\E(x_{k+1})-\E_\infty|^\theta
        \le cL\Delta_k.
\]
Using the concavity of \(t\mapsto t^{1-\theta}\) on positive \(t\), one obtains the standard telescoping estimate
\[
        \Delta_k
        \le
        C\left[
        (\E(x_k)-\E_\infty)^{1-\theta}
        -
        (\E(x_{k+1})-\E_\infty)^{1-\theta}
        \right]
\]
after increasing \(k_1\) if necessary. Summing over \(k\) gives finite length. Hence the tail is Cauchy and converges to \(x_\ast\). Finally, \eqref{eq:relative-error} and \(\Delta_k\to0\) imply \(\nabla\E(x_\ast)=0\).
\end{proof}

\begin{corollary}[Regular-zero convergence]\label{cor:regular-zero}
If the limit point \(x_\ast\) in \cref{thm:kl} satisfies \(G(x_\ast)=0\) and \(DG(x_\ast)\) is invertible, then the Pandrosion sequence converges to the regular zero \(x_\ast\).
\end{corollary}

\begin{proof}
The conclusion is immediate from \cref{thm:kl}. Invertibility identifies the zero as isolated and regular.
\end{proof}

\section{Stable finite-slope steps supply decrease}

\begin{definition}[Stable finite-slope direction]
At \(x\), a finite-slope pair \((x,b)\) is \((K,\mu)\)-stable if \(\Q=\Q_G(x,b)\) is invertible,
\[
        \norm{\Q^{-1}}\le K,
\]
and the direction \(\delta=-\Q^{-1}G(x)\) satisfies
\[
        \langle \nabla\E(x),\delta\rangle\le -\mu\norm{G(x)}^2.
\]
\end{definition}

\begin{proposition}[Line-search sufficient decrease]\label{prop:line-search}
Assume \(\nabla\E\) is Lipschitz with constant \(L_E\) on the relevant level set. If \((x,b)\) is \((K,\mu)\)-stable and \(\delta=-\Q^{-1}G(x)\), then every
\[
        0<\lambda\le \frac{\mu\norm{G(x)}^2}{L_E\norm{\delta}^2}
\]
satisfies
\[
        \E(x+\lambda\delta)
        \le
        \E(x)-\frac{\mu}{2}\lambda\norm{G(x)}^2.
\]
If, in addition, the step sizes satisfy \(\lambda\norm{G(x)}^2\ge c_0\norm{\lambda\delta}^2\), then the accepted steps satisfy sufficient decrease in the form \eqref{eq:suff-decrease}.
\end{proposition}

\begin{proof}
The descent lemma gives
\[
        \E(x+\lambda\delta)
        \le
        \E(x)+\lambda\langle\nabla\E(x),\delta\rangle
        +\frac{L_E}{2}\lambda^2\norm{\delta}^2.
\]
The stability inequality bounds the linear term by \(-\lambda\mu\norm{G(x)}^2\), and the stated step-size bound absorbs half of it. The final statement follows by comparing the residual decrease with \(\norm{\lambda\delta}^2\).
\end{proof}

\section{Gated Halley corrections}

The cubic finite-slope Halley layer should not be accepted unconditionally in a global convergence proof. It must be gated.

\begin{definition}[Descent-compatible gate]
Let \(\delta_k\) be the raw stable Pandrosion direction and \(h_k\) a higher-order correction. A gate is descent-compatible if the accepted correction \(e_k=\omega_k h_k\) satisfies
\[
        \norm{e_k}\le \eta\norm{\delta_k}
\]
with \(\eta>0\) small enough that the line search still accepts a sufficient-decrease step.
\end{definition}

\begin{proposition}[Gated corrections preserve admissibility]\label{prop:gated}
Suppose the raw Pandrosion steps satisfy sufficient decrease and the relative-error condition on a bounded level set. If the gated corrections satisfy \(\norm{e_k}\le\eta\norm{\delta_k}\) with \(\eta\) sufficiently small and the line search is applied to the corrected direction \(\delta_k+e_k\), then the corrected sequence is still Lojasiewicz-admissible, possibly with different constants.
\end{proposition}

\begin{proof}
For small \(\eta\), the corrected direction remains within a fixed cone around the raw descent direction:
\[
        \norm{\delta_k+e_k}\asymp \norm{\delta_k}.
\]
The descent lemma therefore gives the same Armijo-type decrease after reducing the accepted step length by a uniform factor if needed. The relative-error bound is stable under this perturbation because the step norms are comparable and all derivatives are bounded on the level set.
\end{proof}

\section{Pandrosion convergence theorem}

\begin{theorem}[Conditional analytic convergence of gated Pandrosion]\label{thm:pand-convergence}
Let \(G\) be analytic and let a Pandrosion engine generate a bounded sequence using stable finite-slope pairs, line search, and descent-compatible gated corrections. Assume the resulting accepted steps satisfy the relative-error estimate \eqref{eq:relative-error}. Then every cluster point attracts the whole tail of the sequence, and the sequence converges to a critical point of \(\E=\frac12\norm{G}^2\). If the cluster point is a regular zero, the sequence converges to that zero.
\end{theorem}

\begin{proof}
Stable finite-slope pairs and line search give sufficient decrease by \cref{prop:line-search}. Gated corrections preserve admissibility by \cref{prop:gated}. Analyticity gives the Lojasiewicz gradient inequality near cluster points. Apply \cref{thm:kl} and \cref{cor:regular-zero}.
\end{proof}

\section{Global conjecture}

\begin{conjecture}[Regular-attractor Pandrosion conjecture]
For analytic systems whose regular zeros are isolated, a dynamic Pandrosion engine combining stable probe selection, equal-value pole avoidance, \(m\)-flat atlas normalization, and descent-compatible finite-slope Halley gating avoids nonzero critical attractors for an open dense set of initial conditions.
\end{conjecture}

\begin{remark}
This conjecture is stronger than the theorem. The theorem says that admissible trajectories converge. The conjecture says that the full engine can make the regular-zero alternative typical by avoiding singular probes, bad chart regions, and nonzero stationary traps.
\end{remark}

\section{Conclusion}

The Lojasiewicz viewpoint gives the convergence target for the Pandrosion programme. Geometry supplies stable charts and probes; line search supplies Lyapunov decrease; gates keep Halley corrections from destroying descent; analyticity converts these conditions into finite-length convergence. The theorem is conditional, but its conditions are exactly the design rules already isolated in the Pandrosion papers.

\begin{thebibliography}{9}
\bibitem{halley}
I. Besevic, \emph{Cubic Pandrosion by Finite-Slope Halley Acceleration}, manuscript, May 2026.
\bibitem{variational}
I. Besevic, \emph{Variational Principles for Finite-Slope Pandrosion Flows}, manuscript, May 2026.
\bibitem{descent}
I. Besevic, \emph{Geometric Descent Theorems for Pandrosion Flows}, manuscript, May 2026.
\bibitem{atlas}
I. Besevic, \emph{Uniform \(m\)-Flat Pandrosion Atlases}, manuscript, May 2026.
\bibitem{loja}
S. Lojasiewicz, \emph{Une propriete topologique des sous-ensembles analytiques reels}, 1963.
\end{thebibliography}

\end{document}
